\import{set/suc.tex}
\import{set.tex}
\import{order/order.tex}
\import{ordinal.tex}

\section{Natural numbers}

\begin{abbreviation}\label{num_inductive_set}
    $A$ is an inductive set iff $\emptyset\in A$ and for all $a\in A$ we have $\suc{a}\in A$.
\end{abbreviation}

\begin{axiom}\label{num_naturals_inductive_set}
    $\naturals$ is an inductive set.
\end{axiom}

\begin{axiom}\label{num_naturals_smallest_inductive_set}
    Let $A$ be an inductive set.
    Then $\naturals\subseteq A$.
\end{axiom}

\begin{abbreviation}\label{num_naturalnumber}
    $n$ is a natural number iff $n\in \naturals$.
\end{abbreviation}

\begin{lemma}\label{num_emptyset_in_naturals}
    $\emptyset\in\naturals$.
\end{lemma}

\begin{signature}\label{num_addition_is_set}
    $x+y$ is a set.
\end{signature}

\begin{axiom}\label{num_addition_on_naturals}
    If $x$ and $y$ are natural numbers, then $x + y$ is a natural number.
\end{axiom}

\begin{abbreviation}\label{num_zero_is_emptyset}
    $\zero = \emptyset$.
\end{abbreviation}

\begin{axiom}\label{num_addition_axiom_1}
    For all $x \in \naturals$ $x + \zero = \zero + x = x$.
\end{axiom}

\begin{axiom}\label{num_addition_axiom_2}
    For all $x, y \in \naturals$ $x + \suc{y} = \suc{x} + y = \suc{x+y}$.
\end{axiom}

\begin{lemma}\label{num_naturals_is_equal_to_two_times_naturals}
    $\{x+y \mid x \in \naturals, y \in \naturals \} = \naturals$.
\end{lemma}

\begin{lemma}\label{naturals_addition_is_closed}
  Suppose $x$ and $y$ are natural numbers.
  Then $x+y$ is a natural number.
\end{lemma}

\begin{lemma}\label{naturals_addition_cancellation_left}
  For all $x, y, z \in \naturals$ such that $x+y = x+z$ we have $y = z$.
\end{lemma}
\begin{proof}
  Let $S = \{p \in \naturals \mid \text{for all $y,z \in \naturals$ we have if $p + y = p + z$ then $y = z$} \}$.
  For all $y, z \in \naturals$ we have if $\zero + y = \zero + z$ then $y = z$.
  For all $x \in \naturals$ such that $x \in S$, we have for all $y, z \in \naturals$ if $\suc{x} + y = \suc{x} + z$ then $y = z$.
\end{proof}

\begin{lemma}\label{naturals_addition_cancellation_right}
  For all $x, y, z \in \naturals$ such that $y+x = z+x$ we have $y = z$.
\end{lemma}
\begin{proof}
  Let $S = \{p \in \naturals \mid \text{for all $y,z \in \naturals$ we have if $y+p = z+p$ then $y = z$} \}$.
  For all $y, z \in \naturals$ we have if $y + \zero = z + \zero$ then $y = z$.
  For all $x \in \naturals$ such that $x \in S$, we have for all $y, z \in \naturals$ if $y + \suc{x} = z + \suc{x}$ then $y = z$.
\end{proof}

\begin{lemma}\label{naturals_addition_commutative}
  For all $x, y \in \naturals$ we have $x+y = y+x$.
\end{lemma}
\begin{proof}
  Let $S = \{p \in \naturals \mid \text{for all $y \in \naturals$ we have $p+y = y+p$} \}$.
  For all $y \in \naturals$ we have $\zero + y = y + \zero$.
  For all $x \in \naturals$ such that $x \in S$, we have for all $y \in \naturals$ $\suc{x} + y = y + \suc{x}$.
\end{proof}

\begin{lemma}\label{naturals_associative}
  For all $x, y, z \in \naturals$ $x + (y + z) = (x + y) + z$.
\end{lemma}
\begin{proof}
  Let $S = \{p \in \naturals \mid \text{for all $y,z \in \naturals$ we have $p + (y + z) = (p + y) + z$} \}$.
  For all $y, z \in \naturals$ we have $\zero + (y + z) = (\zero + y) + z$.
  For all $x \in \naturals$ such that $x \in S$, we have for all $y, z \in \naturals$ $\suc{x} + (y + z) = (\suc{x} + y) + z$.
\end{proof}

\begin{definition}\label{natural_order}
  $\le = \{(n,m) \mid n \in \naturals, m \in \naturals \mid \text{there exists a natural number $k$ such that $n + k = m$}\}$.
\end{definition}

\begin{lemma}
  \label{naturals_are_ordered_simpler}
  $\le$ is an order on $\naturals$.
\end{lemma}
\begin{proof}
  $\le$ is a binary relation on $\naturals$.

  We show for all $x, y \in \naturals$ such that $x \mathrel{\le} y \mathrel{\le} x$ we have $x = y$.
  \begin{subproof}
    Fix $x, y \in \naturals$.
    Assume $x \mathrel{\le} y \mathrel{\le} x$.
    There exist $a, b, k \in \naturals$ such that $(x,y) = (a,b)$ and $a + k = b$.
    There exist $a, b, k' \in \naturals$ such that $(y,x) = (a,b)$ and $a + k' = b$.
  \end{subproof}
  Therefore $\le$ is antisymmetric.

  For all $x \in \naturals$ we have $x \mathrel{\le} x$.

  We show for all $x, y, z$ such that $x \mathrel{\le} y \mathrel{\le} z$ we have $x \mathrel{\le} z$.
  \begin{subproof}
    Fix $x, y, z$.
    Assume $x \mathrel{\le} y \mathrel{\le} z$.
    Take $k \in \naturals$ such that $x + k = y$.
    Take $k' \in \naturals$ such that $y + k' = z$.
    Then $z = y + k' = (x + k) + k' = x + (k + k')$.
    Therefore $x \mathrel{\le} z$.
  \end{subproof}
\end{proof}
